\chapter{Boolean Analysis --- Threshold Functions --- Bit Flip Sets}
%%% generated by the blueprint pipeline from TCSlib.BooleanAnalysis.ThresholdFunctions.BitFlipSets
\section{Overview}
Setwise bit flips act diagonally on Walsh characters and Fourier coefficients. Their disagreement formula and random-coloring identity quantify simultaneous coordinate flips.

\section{Declarations}

\begin{definition}[Flipping a set of coordinates]
\lean{BooleanAnalysis.ThresholdFunctions.flipSet}
\leanok
\uses{BooleanAnalysis.BoolCube}
For $n\ge0$, $x\in\{0,1\}^n$, and $S\subseteq[n]$, the setwise flip $x^S$ is defined by $(x^S)_i=1-x_i$ if $i\in S$ and $(x^S)_i=x_i$ otherwise.
\end{definition}

\begin{definition}[Coordinatewise addition modulo two]
\lean{BooleanAnalysis.ThresholdFunctions.xorVec}
\leanok
\uses{BooleanAnalysis.BoolCube}
Let $n\ge0$ and $z,c\in\{0,1\}^n$. Their exclusive-or vector $z\oplus c$ has coordinate $(z\oplus c)_i=z_i+c_i\pmod2$.
\end{definition}

\begin{definition}[A color class of a coordinate coloring]
\lean{BooleanAnalysis.ThresholdFunctions.fiber}
\leanok
For integers $n,m\ge0$, a coloring $\pi\colon[n]\to[m]$, and a color $j\in[m]$, its color class is the finite set $F_j(\pi)=\{i\in[n]:\pi(i)=j\}$.
\end{definition}

\begin{lemma}[Membership in a color class]
\lean{BooleanAnalysis.ThresholdFunctions.mem_fiber}
\leanok
\uses{BooleanAnalysis.ThresholdFunctions.fiber}
\difficulty{1}
Let $n,m\ge0$, $\pi\colon[n]\to[m]$, $j\in[m]$, and $i\in[n]$. For $F_j(\pi)=\{u\in[n]:\pi(u)=j\}$, one has $i\in F_j(\pi)$ if and only if $\pi(i)=j$.
\end{lemma}

\begin{lemma}[Flipping all bits gives the antipodal vector]
\lean{BooleanAnalysis.ThresholdFunctions.flipSet_univ}
\leanok
\uses{BooleanAnalysis.BoolCube, BooleanAnalysis.ThresholdFunctions.flipSet}
\difficulty{3}
For $n\ge0$ and $x\in\{0,1\}^n$, flipping all coordinates gives the vector $\bar x$ with $\bar x_i=1-x_i$ for every $i\in[n]$.
\end{lemma}

\begin{lemma}[A fixed setwise flip is an involution]
\lean{BooleanAnalysis.ThresholdFunctions.flipSet_flipSet}
\leanok
\uses{BooleanAnalysis.BoolCube, BooleanAnalysis.ThresholdFunctions.flipSet}
\difficulty{4}
Let $n\ge0$, $x\in\{0,1\}^n$, and $S\subseteq[n]$, and write $x^S$ for the vector formed by flipping exactly the bits in $S$. Then $(x^S)^S=x$.
\end{lemma}

\begin{lemma}[A fixed exclusive-or translation is an involution]
\lean{BooleanAnalysis.ThresholdFunctions.xorVec_xorVec}
\leanok
\uses{BooleanAnalysis.BoolCube, BooleanAnalysis.ThresholdFunctions.xorVec}
\difficulty{3}
For $n\ge0$ and $z,c\in\{0,1\}^n$, let $\oplus$ denote coordinatewise addition modulo two. Then $(z\oplus c)\oplus c=z$.
\end{lemma}

\begin{lemma}[Setwise flips preserve sums over the cube]
\lean{BooleanAnalysis.ThresholdFunctions.sum_comp_flipSet}
\leanok
\uses{BooleanAnalysis.BoolCube, BooleanAnalysis.ThresholdFunctions.flipSet, BooleanAnalysis.ThresholdFunctions.flipSet_flipSet}
\difficulty{1}
Let $n\ge0$, $F\colon\{0,1\}^n\to\mathbb R$, and $S\subseteq[n]$; write $x^S$ for the vector formed by flipping the coordinates in $S$. Then $\sum_{x\in\{0,1\}^n}F(x^S)=\sum_{x\in\{0,1\}^n}F(x)$.
\end{lemma}

\begin{lemma}[Exclusive-or translations preserve cube sums]
\lean{BooleanAnalysis.ThresholdFunctions.sum_comp_xorVec}
\leanok
\uses{BooleanAnalysis.BoolCube, BooleanAnalysis.ThresholdFunctions.xorVec, BooleanAnalysis.ThresholdFunctions.xorVec_xorVec}
\difficulty{1}
Let $n\ge0$, $F\colon\{0,1\}^n\to\mathbb R$, and $c\in\{0,1\}^n$, with $\oplus$ denoting coordinatewise addition modulo two. Then $\sum_{z\in\{0,1\}^n}F(z\oplus c)=\sum_{x\in\{0,1\}^n}F(x)$.
\end{lemma}

\begin{lemma}[Walsh character under a setwise flip]
\lean{BooleanAnalysis.ThresholdFunctions.chiS_flipSet}
\leanok
\uses{BooleanAnalysis.BoolCube, BooleanAnalysis.ThresholdFunctions.flipSet, BooleanAnalysis.boolToSign, BooleanAnalysis.chiS}
\difficulty{4}
Let $n\ge0$ and equip $\{0,1\}^n$ with its uniform probability measure; write $\mathbb E$ for uniform expectation. For $S\subseteq[n]$, put $\chi_S(x)=\prod_{i\in S}(-1)^{x_i}$ and $\widehat h(S)=\mathbb E[h\chi_S]$ for real cube functions $h$. For $S\subseteq[n]$, let $x^S$ be the bit vector obtained by flipping exactly the coordinates in $S$. For all $S,T\subseteq[n]$ and $x\in\{0,1\}^n$, $\chi_T(x^S)=(-1)^{|T\cap S|}\chi_T(x)$.
\end{lemma}

\begin{lemma}[Fourier coefficients under setwise flips]
\lean{BooleanAnalysis.ThresholdFunctions.fourierCoeff_comp_flipSet}
\leanok
\uses{BooleanAnalysis.BooleanFunc, BooleanAnalysis.ThresholdFunctions.chiS_flipSet, BooleanAnalysis.ThresholdFunctions.flipSet, BooleanAnalysis.ThresholdFunctions.flipSet_flipSet, BooleanAnalysis.ThresholdFunctions.sum_comp_flipSet, BooleanAnalysis.chiS, BooleanAnalysis.expect, BooleanAnalysis.fourierCoeff, BooleanAnalysis.innerProduct}
\difficulty{3}
Let $n\ge0$ and equip $\{0,1\}^n$ with its uniform probability measure; write $\mathbb E$ for uniform expectation. For $S\subseteq[n]$, put $\chi_S(x)=\prod_{i\in S}(-1)^{x_i}$ and $\widehat h(S)=\mathbb E[h\chi_S]$ for real cube functions $h$. For $S\subseteq[n]$, let $x^S$ be the bit vector obtained by flipping exactly the coordinates in $S$. For every real cube function $f$ and sets $S,T\subseteq[n]$, $\widehat{(x\mapsto f(x^S))}(T)=(-1)^{|T\cap S|}\widehat f(T)$.
\end{lemma}

\begin{lemma}[Fourier formula for setwise-flip correlation]
\lean{BooleanAnalysis.ThresholdFunctions.expect_mul_flipSet}
\leanok
\uses{BooleanAnalysis.BooleanFunc, BooleanAnalysis.ThresholdFunctions.flipSet, BooleanAnalysis.ThresholdFunctions.fourierCoeff_comp_flipSet, BooleanAnalysis.expect, BooleanAnalysis.fourierCoeff, BooleanAnalysis.innerProduct, BooleanAnalysis.plancherel}
\difficulty{3}
Let $n\ge0$ and equip $\{0,1\}^n$ with its uniform probability measure; write $\mathbb E$ for uniform expectation. For $S\subseteq[n]$, put $\chi_S(x)=\prod_{i\in S}(-1)^{x_i}$ and $\widehat h(S)=\mathbb E[h\chi_S]$ for real cube functions $h$. For $S\subseteq[n]$, let $x^S$ be the bit vector obtained by flipping exactly the coordinates in $S$. For every real cube function $f$ and $S\subseteq[n]$, $\mathbb E[f(x)f(x^S)]=\sum_{T\subseteq[n]}(-1)^{|T\cap S|}\widehat f(T)^2$.
\end{lemma}

\begin{definition}[Quadratic disagreement under a setwise flip]
\lean{BooleanAnalysis.ThresholdFunctions.flipDisagreement}
\leanok
\uses{BooleanAnalysis.BooleanFunc, BooleanAnalysis.ThresholdFunctions.flipSet, BooleanAnalysis.expect}
Let $n\ge0$ and equip $\{0,1\}^n$ with its uniform probability measure; write $\mathbb E$ for uniform expectation. For $S\subseteq[n]$, let $x^S$ be the bit vector obtained by flipping exactly the coordinates in $S$. For a real cube function $f$ and $S\subseteq[n]$, define $d_S(f)=\tfrac14\mathbb E[(f(x)-f(x^S))^2]$. For sign-valued $f$, this equals $\Pr[f(x)\ne f(x^S)]$.
\end{definition}

\begin{lemma}[Fourier formula for flip disagreement]
\lean{BooleanAnalysis.ThresholdFunctions.flipDisagreement_eq}
\leanok
\uses{BooleanAnalysis.BooleanFunc, BooleanAnalysis.ThresholdFunctions.expect_add, BooleanAnalysis.ThresholdFunctions.expect_const, BooleanAnalysis.ThresholdFunctions.expect_const_mul, BooleanAnalysis.ThresholdFunctions.expect_mul_flipSet, BooleanAnalysis.ThresholdFunctions.flipDisagreement, BooleanAnalysis.ThresholdFunctions.flipSet, BooleanAnalysis.fourierCoeff, BooleanAnalysis.isPmOne}
\difficulty{4}
Let $n\ge0$ and equip $\{0,1\}^n$ with its uniform probability measure; write $\mathbb E$ for uniform expectation. For $S\subseteq[n]$, put $\chi_S(x)=\prod_{i\in S}(-1)^{x_i}$ and $\widehat h(S)=\mathbb E[h\chi_S]$ for real cube functions $h$. For $S\subseteq[n]$, let $x^S$ be the bit vector obtained by flipping exactly the coordinates in $S$. If $f\colon\{0,1\}^n\to\{-1,1\}$, then $\Pr[f(x)\ne f(x^S)]=(1-\sum_{T\subseteq[n]}(-1)^{|T\cap S|}\widehat f(T)^2)/2$ for each $S\subseteq[n]$; equivalently, this is $\tfrac14\mathbb E[(f(x)-f(x^S))^2]$.
\end{lemma}

\begin{lemma}[Color-class intersection parity as a product]
\lean{BooleanAnalysis.ThresholdFunctions.neg_one_pow_card_inter_fiber}
\leanok
\uses{BooleanAnalysis.ThresholdFunctions.fiber}
\difficulty{4}
Let $n,m\ge0$, $j_0\in[m]$, $T\subseteq[n]$, and $\pi\colon[n]\to[m]$, and put $F_{j_0}(\pi)=\{i:\pi(i)=j_0\}$. Then $(-1)^{|T\cap F_{j_0}(\pi)|}=\prod_{i=1}^n b_i$, where $b_i=-1$ if $i\in T$ and $\pi(i)=j_0$, and $b_i=1$ otherwise.
\end{lemma}

\begin{lemma}[Parity sum over coordinate colorings]
\lean{BooleanAnalysis.ThresholdFunctions.sum_pi_sign}
\leanok
\uses{BooleanAnalysis.ThresholdFunctions.fiber, BooleanAnalysis.ThresholdFunctions.neg_one_pow_card_inter_fiber}
\difficulty{5}
Let $n\ge0$, $m>0$ be integers, $j_0\in[m]$, and $T\subseteq[n]$. For $F_{j_0}(\pi)=\{i\in[n]:\pi(i)=j_0\}$, the sum over all colorings satisfies $\sum_{\pi:[n]\to[m]}(-1)^{|T\cap F_{j_0}(\pi)|}=m^n(1-2/m)^{|T|}$.
\end{lemma}
